\section{Evaluation}\label{sec:eval}

\begin{table}[t]
\caption{The ten benchmark algorithms, with the asymptotic complexity reduction delivered by a single \code{@Memoize} annotation.}
\label{tab:algorithms}
\small
\centering
\begin{tabularx}{\columnwidth}{@{}l l l@{}}
\toprule
\textbf{Algorithm} & \textbf{Naive} & \textbf{Memoized} \\
\midrule
Fibonacci                 & $O(\varphi^n)$          & $O(n)$ \\
Ackermann                 & $\gg$ exponential       & O(reachable lattice) \\
Catalan                   & $O(C_n)$                & $O(n^2)$ \\
Binomial (Pascal)         & $O(2^n)$                & $O(n k)$ \\
Integer partition         & exponential             & $O(n k)$ \\
Edit distance             & $O(3^{n+m})$            & $O(n m)$ \\
Longest common subseq.    & $O(2^{n+m})$            & $O(n m)$ \\
Matrix chain              & $O(2^n)$                & $O(n^3)$ \\
0/1 knapsack              & $O(2^n)$                & $O(n W)$ \\
Coin change (ways)        & exponential             & $O(n A)$ \\
\bottomrule
\end{tabularx}
\end{table}


\paragraph{Scope of the empirical claims.} Three of the design goals---usability, automation-friendliness, and cross-platform reach---are qualitative contributions argued on architectural grounds in Section~\ref{sec:tool} (annotation surface, build-time rewrite, zero-dependency runtime). The empirical evaluation below is deliberately scoped to the remaining, measurable claim: the performance cost of the abstraction itself, in both its enabled and disabled regimes.

To demonstrate that the abstraction does not sacrifice the performance it is meant to deliver, we release alongside the library an open Jetpack Benchmark~\cite{jetpackbenchmark} harness over ten compute-intensive, overlapping-subproblem algorithms: naive Fibonacci, Ackermann, Catalan, Pascal-recurrence binomial, integer partition, Levenshtein edit distance, longest common subsequence, matrix-chain multiplication, 0/1 knapsack, and coin change. Each algorithm exists in a baseline un-memoized form and in three memoized subclasses, one per eviction policy; subclasses \code{@Override} every method with a one-line \code{super} call carrying the \code{@Memoize} annotation, so that the parent's recursive \code{invokevirtual} self-calls dispatch back through the memoized subclass and cascade the cache through the entire recursion tree without source duplication. Because both baseline and memoized forms already dispatch through \code{invokevirtual}, this is a fair comparison; readers should nevertheless note that the memoized path incurs one extra call-frame through \code{super}, so the reported speedups slightly \emph{under}-report the raw algorithmic speedup attributable to memoization alone.

Table~\ref{tab:algorithms} reports the asymptotic complexity reduction delivered by a single \code{@Memoize} annotation. Four benchmark suites operationalise four research questions. \textbf{RQ1} pairs each baseline with its \lru-memoized variant to measure per-algorithm speedup. \textbf{RQ2} compares \lru, \fifo, and \lfu on workloads that fit within the cache (pure hot-path cost) and one that forces eviction under a skewed distribution. \textbf{RQ3} sweeps Fibonacci across four input distributions (HOT, WARM, COLD, VERY\_COLD) to locate the hit-rate regime at which \code{autoMonitor} should intervene. \textbf{RQ4} quantifies the library's per-call wrapper overhead in isolation through a dedicated \code{OverheadBenchmark}, described below. The harness builds against \code{testBuildType~=~release} with R8 enabled, so the reported medians reflect production inlining, devirtualisation, and \dce; sink values are \code{volatile}-assigned to prevent full-mode \dce from eliminating the measured calls.

\paragraph{Isolating library overhead (RQ4).}
\code{OverheadBenchmark} drives a fixed workload of 100 calls per iteration through a synthetic, non-recursive \code{expensiveCompute(int)} whose per-call compute cost is constant (a 500-iteration mixing loop independent of its input). Each iteration pre-populates a fresh \code{OverheadLibLru} with 8 hot keys during a \code{pauseTiming()} window, then runs a deterministically shuffled mix of hot and unique cold keys selected so the observed steady-state hit rate matches the target exactly. \code{maxSize}~=~128 is large enough that no cold key can evict a hot key within the iteration, decoupling the measurement from eviction effects. Four target rates are exercised---75\%, 50\%, 25\%, and 0\%---against a single baseline row that runs the same workload length through the un-memoized \code{OverheadLib}. The per-call overhead attributable to the library is $(\text{memoized}_{ns} - \text{baseline}_{ns})/100$: negative values indicate that compute savings on hits outweigh the wrapper cost; positive values the opposite. Because \code{expensiveCompute} has constant cost and no recursion, the measured delta is attributable entirely to the library's wrapper, not to any algorithmic confound.

\begin{table}[t]
\caption{\code{OverheadBenchmark} results: per-call wrapper overhead of \code{@Memoize} at four controlled hit rates. The synthetic target (\code{expensiveCompute}) has a constant per-call compute cost, so the delta is attributable entirely to the library.}
\label{tab:results_overhead}
\small
\centering
\begin{tabularx}{\columnwidth}{@{}l r r r@{}}
\toprule
\textbf{Hit rate} & \textbf{Baseline (ns)} & \textbf{Memoized (ns)} & \textbf{$\Delta$/call (ns)} \\
\midrule
75\% & \pending & \pending & \pending \\
50\% & \pending & \pending & \pending \\
25\% & \pending & \pending & \pending \\
0\%  & \pending & \pending & \pending \\
\bottomrule
\end{tabularx}
\end{table}


Table~\ref{tab:results_overhead} presents the results. The hit rate at which the per-call overhead crosses zero is the empirical break-even point: above it, \code{@Memoize} is strictly beneficial; below it, the library is costing more time than it saves. The same break-even doubles as a data-driven lower bound for \code{autoMonitor}: configuring \code{@Memoize} with \code{autoMonitor~=~true} and \code{minHitRate} set to $b + \varepsilon$, where $b$ is the observed break-even, guarantees that any method whose runtime hit rate falls below $b$ is automatically disabled after one monitor window. Because the target's compute cost is fixed in this benchmark, the break-even hit rate scales predictably with real targets---cheaper compute raises it, more expensive compute lowers it---which lets developers pick a conservative \code{minHitRate} even when they only measure one reference point.

The expected shape of the results across all four suites follows directly from algorithmic complexity and the implementation details of each policy. For \RQ{1}, every row is expected to deliver several orders of magnitude of speedup, because memoization converts exponential recurrences into polynomial cached computations: a single \code{@Memoize} annotation flips naive Fibonacci from $O(\varphi^n)$ to $O(n)$, binomial from $O(2^n)$ to $O(nk)$, and matrix-chain from $O(2^n)$ to $O(n^3)$. For \RQ{2}, the no-eviction rows are expected to order as \fifo~$<$~\lru~$<$~\lfu, reflecting the per-operation tax of each policy (no re-linking, one re-linking, frequency-bucket bookkeeping); under eviction pressure with a skewed distribution, \lfu is expected to out-hit \lru. For \RQ{3}, the break-even is expected to lie in the VERY\_COLD regime (hit rate $\approx$~0.05), well below the default \code{autoMonitor} threshold of~0.3. Absolute Pixel~8 and Pixel~10 medians populate the result tables in the artefact's replication package as the device campaign is executed.

\paragraph{Threats to validity.} Four caveats qualify the results. (a)~Numeric cells are collected on a single device class (Pixel~8, Pixel~10); behaviour on other \art builds or chipsets is not attested. (b)~\code{OverheadBenchmark}'s target is synthetic with constant per-call cost, which isolates the wrapper but does not model compute-cost variance across arguments. (c)~R8 inlining, devirtualisation, and \dce behaviour may shift across releases. (d)~Result cells marked \pending{} await the full device campaign; until replaced, only the shape predictions---not the absolute numbers---should be relied upon.
